
\bigskip

\textbf{Partie 1}

\medskip

Soit $g$ la fonction définie pour tout réel $x\in \; ]0~;~+ \infty[$ par:
$g(x) = \dfrac{2 \ln x}{ x}.$

%\smallskip

\begin{enumerate}
\item On note $g'$ la dérivée de $g$.
On a: $g'(x)=\dfrac{2\dfrac{1}{x}\times x - 2\ln x \times 1}{x^2}
= \dfrac{2-2\ln x}{x^2}$

\item On dispose de ce tableau de variations de la fonction $g$ sur l'intervalle $]0~;~+ \infty[$ : 

\begin{center}
{\renewcommand{\arraystretch}{1.3}
\psset{nodesep=3pt,arrowsize=2pt 3}%  paramètres
\def\esp{\hspace*{1.5cm}}% pour modifier la largeur du tableau
\def\hauteur{20pt}% mettre au moins 20pt pour augmenter la hauteur
$\begin{array}{|c|l*4{c}|}
\hline
x & 0  & \esp\esp & \e & \esp & +\infty \\ 
%\hline
%f'(x) &  &   \pmb{+} & \vline\hspace{-2.7pt}0 & \pmb{-} & \\ 
\hline
 &\vline\;\vline &  &   \Rnode{max}{\frac{2}{\e}}  &  &   \\  
\text{Variations} \atop \text{de }g & \vline\;\vline &     &  &  &  \rule{0pt}{\hauteur} \\ 
 &\vline\;\vline \Rnode{min1}{~-\infty} &   &  &  &   \Rnode{min2}{0} \rule{0pt}{\hauteur}    
 \ncline{->}{min1}{max} 
 \ncline{->}{max}{min2} 
 \rput*(-5,0.915){\Rnode{zero}{0}}
\rput(-5,2.55){\Rnode{alpha}{1}}
%\ncline[linestyle=dotted, linecolor=blue]{alpha}{zero}
%\rput*(-1.3,0.65){\Rnode{zero2}{\red 0}}
%\rput(-1.3,2.4){\Rnode{beta}{\red \beta}}
%\ncline[linestyle=dotted, linecolor=red]{beta}{zero2}
 \\ 
\hline
\end{array} $
}
\end{center}	

%\smallskip
%
%Justifier les informations suivantes lues dans ce tableau:

	\begin{enumerate}
		\item La valeur $\dfrac{2}{\text{e}}$ est l'image de e par $f$: $f(\e)=\dfrac{2\ln\e}{\e} = \dfrac{2}{\e}$.
		\item %les variations de la fonction $g$ sur son ensemble de définition ;
$g'(x)= \dfrac{2-2\ln x}{x^2}$ donc $g'(x)$ est du signe de $2-2\ln x=2\left (1-\ln x\right )$.
\begin{list}{\textbullet}{}
\item Sur $]0\;;\e[$, $\ln x < 1$ donc $1-\ln x>0$ donc $g'(x)>0$; la fonction $g$ est strictement croissante sur cet intervalle.
\item Sur $]\e\;; +\infty[$, $\ln x >1$ donc $1-\ln x<0$ donc $g'(x)<0$; la fonction $g$ est strictement décroissante sur cet intervalle.
\end{list}		

		\item On justifie les limites de la fonction $g$ aux bornes de son ensemble de définition.
\begin{list}{\textbullet}{}
\item $\ds\lim_{x\to 0\atop x>0} 2\ln x = -\infty$ et $\ds\lim_{x\to 0\atop x>0} \dfrac{1}{x}=+\infty$. \\
Donc, par produit de limites, $\ds\lim_{x\to 0\atop x>0} g(x)= \ds\lim_{x\to 0\atop x>0} 2\ln x\times \dfrac{1}{x}=-\infty$.
\item D'après le cours, $\ds\lim_{x\to +\infty} \dfrac{\ln x}{x}=0$ donc $\ds\lim_{x\to +\infty} g(x)=0$. 
\end{list}		
		
	\end{enumerate}
	
\item On en déduit le tableau de signes de la fonction $g$ sur l'intervalle $]0~;~+ \infty[$.

\begin{center}
\renewcommand{\arraystretch}{1.2}
\def\esp{\hspace*{1.1cm}}
$\begin{array}{|c |l *{3}{c} r|} 
\hline
x  & 0 & \esp & 1 & \esp\esp  & +\infty \\
\hline
\text{signe de }g(x) & \vline\;\vline  & \pmb{-} &  \vline\hspace{-2.7pt}{0} & \pmb{+} & 0   \\
\hline
\end{array}$
\renewcommand{\arraystretch}{1}
\end{center}

\end{enumerate}

\bigskip

\textbf{Partie 2}

\medskip

Soit $f$ la fonction définie sur l'intervalle $]0~;~ + \infty[$ par 
$f(x) = [\ln (x)]^2.$

 Dans cette partie, chaque étude est effectuée sur l'intervalle $]0~;~+ \infty[$.

%\medskip

\begin{enumerate}
\item %Démontrer que sur l'intervalle $]0~;~+ \infty[$, la fonction $f$ est une primitive de la fonction $g$.
$[(u(x)^2]' = 2u'(x)u(x)$ donc
$f'(x)= 2\times \dfrac{1}{x}\times \ln x = \dfrac{2\ln x}{x}=g(x)$

Donc sur l'intervalle $]0~;~+ \infty[$, la fonction $f$ est une primitive de la fonction $g$.

\item %À l'aide de la partie 1, étudier :
	\begin{enumerate}
		\item On étudie la convexité de la fonction $f$.
		
D'après les questions précédentes, la fonction $g$, dérivée de la fonction $f$, est croissante sur $]0\;; \e[$, donc la fonction $f$ est convexe sur cet intervalle.		
		
De même, la fonction $g$, dérivée de la fonction $f$, est décroissante sur $]\e\;; +\infty[$, donc la fonction $f$ est concave sur cet intervalle.

De plus, la fonction $g$ donc la fonction $f'$, change de sens de variation en $x=\e$, donc la courbe représentant la fonction $f$ admet un point d'inflexion en $x=\e$.
		
		\item On étudie les variations de la fonction $f$ en utilisant le signe de $f'=g$.
		
Sur l'intervalle $]0\;; 1[$, la fonction $g$ est négative donc $f'$ est négative; la fonction $f$ est donc strictement décroissante sur cet intervalle.

Sur l'intervalle $]1\;; +\infty[$, la fonction $g$ est positive donc $f'$ est positive; la fonction $f$ est donc strictement croissante sur cet intervalle.

De plus on peut dire que la fonction $f$ admet un minimum pour $x=1$.
		
\end{enumerate}
\item 
	\begin{enumerate}
		\item Une équation de la tangente à la courbe représentative de $f$ au point d'abscisse e est: $y=f'(\e)\left (x-\e\right )+f(\e)$. On a: $f'(\e) = g(\e)=\dfrac{2}{\e}$; $f(\e)=\left (\ln\e\right )^2=1$

L'équation devient: $y=\dfrac{2}{\e}\left (x-\e\right ) +1$ soit
$y=\dfrac{2}{\e}x -2+1$ c'est-à-dire $y=\dfrac{2}{\e}x-1$.
		
		\item %En déduire que, pour tout réel $x$ dans $]0~;~\text{e}]$: $[\ln (x)]^2 \geqslant \dfrac{2}{\text{e}}x - 1.$
La courbe représentant $f$ admet un point d'inflexion en $x=\e$ donc la tangente en ce point coupe la courbe; à gauche du point, la fonction est convexe donc la courbe est au dessus de cette tangente.

On en déduit que sur  $]0~;~\e]$, on a: $[\ln (x)]^2 \geqslant \dfrac{2}{\e}x - 1.$

	\end{enumerate}
\end{enumerate}

\bigskip

